\section{Basics}

\begin{frame}{Lists and emphasis}
  \begin{itemize}
    \item A plain bullet point
    \item Emphasis with \hl{the highlight macro}
    \item De-emphasis with \muted{the muted macro}
      \begin{itemize}
        \item A nested point
          \begin{itemize}
            \item And a third level
          \end{itemize}
      \end{itemize}
    \item An \alert{alerted} word
  \end{itemize}

  \begin{enumerate}
    \item Numbered lists work too
    \item Second item
  \end{enumerate}
\end{frame}

\begin{frame}{Blocks}
  \begin{block}{Standard block}
    Used for definitions and key statements.
  \end{block}

  \begin{exampleblock}{Example block}
    Used for worked examples.
  \end{exampleblock}

  \begin{alertblock}{Alert block}
    Used for warnings and caveats.
  \end{alertblock}
\end{frame}

\begin{frame}{Two columns and a table}
  \twocol{%
    \begin{itemize}
      \item The \code{\textbackslash twocol} macro
      \item Takes left and right content
      \item Optional width: \code{[0.6]}
    \end{itemize}
  }{%
    \begin{table}
      \centering
      \small
      \begin{tabular}{lrr}
        \toprule
        Method & Accuracy & Time \\
        \midrule
        Baseline & 82.1 & 1.0 \\
        Ours     & \textbf{91.4} & 1.3 \\
        \bottomrule
      \end{tabular}
    \end{table}
  }
\end{frame}

\begin{frame}{Overlays}
  Content can appear step by step:
  \begin{itemize}
    \item<1-> Visible from the first click
    \item<2-> Appears second
    \item<3-> Appears third
  \end{itemize}

  \onslide<3->{\hl{All three are now showing.}}
\end{frame}

\statementframe{A statement slide, for one big idea}
